\chapter{Радиальная устойчивость тетраэдрического вихря}

Существование стационарного профиля ещё не исключает отрицательную моду.
Здесь вычисляется второй вариационный оператор в осесимметричном секторе
трёх связанных возмущений $\delta K$, $\delta a$, $\delta b$.

Второй вариационный функционал включает правильные радиальные меры:
для калибровочного профиля $G/r$, для двух скалярных профилей $Ar$ и
$Br$. Смешанные блоки равны
\begin{equation}
 H_{Ka}=-\frac{2Ca(1-K)}r,qquad
 H_{ab}=r\,\partial_a\partial_bV.
\end{equation}
На $K$ и $a$ наложены однородные условия Дирихле на обоих концах, на
$b$ --- естественное условие Неймана в ядре и Дирихле на бесконечности.

Оператор собран кусочно-линейным методом конечных элементов. При числе
узлов $100,160,240,360,500$ минимальное собственное значение равно
соответственно
\begin{equation}
 1.51051387, 1.51049856, 1.51046835, 1.51044709, 1.51043666.
\end{equation}
Сдвиг между двумя последними сетками равен $1.04\cdot10^{-5}$.

На сетке из пятисот узлов первые восемь собственных значений равны
\begin{equation}
 \{1.51044,1.59558,1.70182,1.86176,
 2.07352,2.33631,2.64961,3.01290\}.
\end{equation}
Все они положительны; нулевых и отрицательных радиальных мод нет.

\begin{theorem}[Осесимметричная радиальная устойчивость]
В трёхпрофильном секторе $K+a+b$ найденный $Z_3$-вихрь является строгим
локальным минимумом. Его радиальный гессиан имеет положительную щель не
меньше $1.5104$ в проверенном пределе сгущения сетки.
\end{theorem}

Этот результат не является полной устойчивостью. Не проверены угловые
гармоники, продольные волновые числа и полный калибровочно фиксированный
оператор всех компонент $Q$, $T$ и $B$. Общая литература по вихрям
показывает, что именно внерадиальные или длинноволновые каналы могут
содержать отрицательные моды; см. \path{arXiv:0712.3589} и
\path{arXiv:2608.10608}.

\begin{protocol}
\begin{enumerate}
 \item радиальный гессиан: \textbf{построен};
 \item отрицательные радиальные моды: \textbf{нет};
 \item нулевые радиальные моды: \textbf{нет};
 \item сеточная сходимость: \textbf{получена};
 \item полная устойчивость: \textbf{открыта};
 \item следующий гейт: \textbf{угловые сектора вихря}.
\end{enumerate}
\end{protocol}

Машинный сертификат сохранён в
\path{s2t_v6_bosonic_defect_q_tetrahedral_vortex_radial_stability_gate_results.json}.